\documentclass[12pt,a4paper]{article}

\usepackage{fontspec}
\setmainfont{Latin Modern Roman}
\newfontfamily{\hindifont}{Nirmala UI}[Script=Devanagari]
\newcommand{\hindi}[1]{{\hindifont #1}}

\usepackage[margin=2.5cm]{geometry}
\usepackage{xcolor}
\usepackage{enumitem}
\usepackage{hyperref}
\usepackage{tcolorbox}
\usepackage{booktabs}
\usepackage{tabularx}
\usepackage{fancyhdr}
\usepackage{titlesec}
\usepackage{setspace}

\tcbuselibrary{skins, breakable}

% Colours
\definecolor{govblue}{HTML}{003366}
\definecolor{govgold}{HTML}{C19A3E}
\definecolor{accent1}{HTML}{1B4F72}
\definecolor{accent2}{HTML}{2E86C1}
\definecolor{accent3}{HTML}{27AE60}
\definecolor{accent4}{HTML}{E67E22}
\definecolor{accent5}{HTML}{8E44AD}
\definecolor{redalert}{HTML}{E53935}

% tcolorbox styles
\newtcolorbox{datebox}[1][]{
  colback=govblue!5, colframe=govblue,
  fonttitle=\bfseries, title={#1},
  boxrule=0.5pt, arc=2mm, breakable
}

\newtcolorbox{actionbox}[1][]{
  colback=accent4!8, colframe=accent4,
  fonttitle=\bfseries, title={#1},
  boxrule=0.5pt, arc=2mm, breakable
}

\newtcolorbox{insightbox}[1][]{
  colback=accent3!8, colframe=accent3,
  fonttitle=\bfseries, title={#1},
  boxrule=0.5pt, arc=2mm, breakable
}

\newtcolorbox{criticalbox}[1][]{
  colback=redalert!8, colframe=redalert,
  fonttitle=\bfseries, title={#1},
  boxrule=0.5pt, arc=2mm, breakable
}

% Header/Footer
\pagestyle{fancy}
\fancyhf{}
\fancyhead[L]{\small\textcolor{govblue}{AIGGPA Internship -- Notes from Manager}}
\fancyhead[R]{\small\textcolor{govblue}{April 2026}}
\fancyfoot[C]{\thepage}
\renewcommand{\headrulewidth}{0.4pt}
\setlength{\headheight}{14pt}

% Section Formatting
\titleformat{\section}
  {\Large\bfseries\color{govblue}}{\thesection}{1em}{}[\titlerule]
\titleformat{\subsection}
  {\large\bfseries\color{accent1}}{\thesubsection}{1em}{}
\titleformat{\subsubsection}
  {\normalsize\bfseries\color{accent2}}{\thesubsubsection}{1em}{}

\setstretch{1.25}
\hypersetup{colorlinks=true, linkcolor=govblue, urlcolor=accent2}

\begin{document}

% --- Title ---
\begin{center}
\vspace*{0.5cm}

{\Huge\bfseries\textcolor{govblue}{Notes from Manager}}

\vspace{0.3cm}

{\Large\textcolor{accent1}{AIGGPA Bhopal -- Summer Internship 2026}}

{\large\textcolor{accent1}{\hindi{एआईजीजीपीए भोपाल -- ग्रीष्मकालीन इंटर्नशिप}}}

\vspace{0.5cm}

\rule{0.6\textwidth}{1pt}

\vspace{0.3cm}

{\large Handwritten field notes transcribed from internship diary}

{\normalsize Covering: 6 April 2026 -- 9 April 2026}

\end{center}

\vspace{1cm}

\tableofcontents
\newpage

% =============================================================================
% DAY 1 -- 6 APRIL 2026
% =============================================================================
\section{Day 1: 6 April 2026 -- Initial Briefing \& Task Assignment}

\begin{datebox}[To-Do List from Manager]
\textbf{Research Tasks Assigned:}
\begin{enumerate}[label=\textbf{\arabic*.}, leftmargin=1.5cm]
  \item[\textbf{K} $\rightarrow$] Work on given topics of study
  \item[\textbf{K} $\rightarrow$] \textbf{HR \& IT in Government Sector}
    \begin{itemize}
      \item Focus areas: women, senior officials, gadgets/devices
    \end{itemize}
  \item[\textbf{S} $\rightarrow$] \textbf{Mission Karmayogi} (\hindi{मिशन कर्मयोगी})
  \item[\textbf{H} $\rightarrow$] \textbf{Wellness of HR} -- Is HR productive?
    \begin{itemize}
      \item Study different aspects of wellness (\textbf{8 aspects})
    \end{itemize}
\end{enumerate}
\end{datebox}

\begin{actionbox}[Immediate Action Item]
\textbf{Monday Sunrise $\rightarrow$ Presentation}

\noindent Research the following for the presentation:
\begin{itemize}
  \item DoPT website and other platforms
  \item CDEC (Centre for Distance \& Online Education)
  \item NITI Aayog
  \item Government Departments and their platforms
\end{itemize}
\end{actionbox}

% ---
\subsection{Key Investigation Areas}

The manager outlined three layers of investigation for understanding technology usage among government employees:

\begin{enumerate}[label=\textbf{Layer \arabic*:}]
  \item \textbf{Government Departments that are not working efficiently} -- or at least, not as efficiently as they can
  \item \textbf{Ask government department employees:}
    \begin{itemize}
      \item[$\rightarrow$] Who \textit{is} using IT tools?
      \item[$\rightarrow$] Who \textit{knows about} IT tools but is \textit{not using} them?
      \item[$\rightarrow$] If HR is using tools -- who does not even \textit{know} about them?
      \item[$\rightarrow$] Or not using at all
    \end{itemize}
\end{enumerate}

\begin{insightbox}[Key Insight]
This is essentially a \textbf{3-tier technology awareness classification}:
\begin{enumerate}
  \item \textbf{Active Users} -- using IT tools
  \item \textbf{Aware Non-Users} -- know about tools but don't use them
  \item \textbf{Unaware} -- don't even know the tools exist
\end{enumerate}
This directly maps to the adoption index methodology in our research framework.
\end{insightbox}

% =============================================================================
\section{iGOT Karmayogi Partners \& Training Programmes}

\subsection{iGOT Partners (Content \& Training Providers)}

\begin{table}[h!]
\centering
\caption{iGOT Karmayogi Partner Institutions}
\begin{tabularx}{\textwidth}{@{}lX@{}}
\toprule
\textbf{Partner} & \textbf{Category} \\
\midrule
LBSNAA (Mussoorie) & Premier Civil Services Training Academy \\
ISTM & Institute of Secretariat Training \& Management \\
SVNPA & Sardar Vallabhbhai National Police Academy \\
ISRO & Indian Space Research Organisation \\
Microsoft & Technology Partner \\
Udemy & Online Learning Platform \\
Karmayogi Bharat (\hindi{कर्मयोगी भारत}) & Nodal Agency for Mission Karmayogi \\
Indian Institute of Remote Sensing, Dehradun & Specialised Technical Training \\
IIT Roorkee & Academic Partner \\
\bottomrule
\end{tabularx}
\end{table}

\subsection{Training Programme Topics (from iGOT/LBSNAA)}

\subsubsection{Technical/Domain Training}
\begin{enumerate}
  \item Application of AI/ML in Agricultural Analytics
  \item Water Management / Natural Resource Management
  \item Blue Economy (\hindi{नीली अर्थव्यवस्था})
\end{enumerate}

\subsubsection{Governance \& Policy Training (via LBSNAA)}
\begin{enumerate}
  \item Driving Government Finance for Developed India
  \item Critical Issues of Governance in PRIs \& Rural Development
  \item Border Economy \& Border Development for National Security
  \item Breaking Barriers: Building Property Strategies for Women's Economic Inclusion
  \item Capacity Building \& Sensitisation Programme on Criminal Law Reform
  \item Digital Transformation of Government for Improved Public Service
\end{enumerate}

\subsubsection{Ethics \& Personal Development}
\begin{enumerate}
  \item Ethics in Public Service (via Centre for Governance)
  \item Building Competencies for Personal Excellence (via Art of Living, Bangalore)
\end{enumerate}

\noindent\textit{Additional reference: LBSNAA, Jangpok}

% =============================================================================
\newpage
\section{Day 4: 9 April 2026 -- Mission Karmayogi Deep Dive}

\begin{datebox}[Mission Karmayogi (\hindi{मिशन कर्मयोगी}) -- Briefing Notes]

\textbf{What is Mission Karmayogi?}

An Indian Government initiative to \textbf{transform} bureaucratic learning from \textbf{RULE-BASED} to \textbf{ROLE-BASED}.

\vspace{0.5em}

\textbf{Two Core Components:}
\begin{enumerate}[label=\textbf{(\arabic*)}]
  \item \textbf{iGOT Karmayogi Platform} \\
    Enables officials to pursue learning \textbf{anytime, anyplace, any device} at their own convenience.
  \item \textbf{National Programme for Civil Service Capacity Building (NPCSCB)}
\end{enumerate}
\end{datebox}

\subsection{6 Pillars of Mission Karmayogi}

\begin{table}[h!]
\centering
\caption{Six Pillars of Mission Karmayogi}
\begin{tabularx}{\textwidth}{@{}clX@{}}
\toprule
\textbf{No.} & \textbf{Pillar} & \textbf{Description} \\
\midrule
1 & iGOT Karmayogi & Large-scale comprehensive digital learning platform \\
2 & e-HRMS & Strategic HR management platform \\
3 & M\&E & Continuous performance analysis \\
4 & Institutional Framework & Data-driven goal setting; oversight by PM HR Council \\
5 & Competence Framework & Shift from Rule $\rightarrow$ Role; new training policies \\
6 & Policy Framework & Focus on outcome-based (end) learning \\
\bottomrule
\end{tabularx}
\end{table}

% =============================================================================
\section{Field Research Observations \& Questions}

\subsection{Critical Finding: Research Gap}

\begin{criticalbox}[No Existing Research Found]
\begin{enumerate}[label=\textbf{(\arabic*)}]
  \item \textbf{No research exists on HR \& IT integration} -- no papers available
  \item As of T4 2026, research on HR: approximately \textbf{22 papers found} (none on HR-IT specifically)
  \item \textbf{No existing dataset in AIGGPA}
  \item \textbf{Do people in AIGGPA use e-Office?} -- Status unclear (needs investigation)
  \item Archiving and digitisation of files in AIGGPA -- needs assessment
\end{enumerate}
\end{criticalbox}

\subsection{Manager's Questions for Employees}

These are the questions the manager wants investigated through primary research:

\begin{table}[h!]
\centering
\caption{Research Questions from Manager}
\begin{tabularx}{\textwidth}{@{}clX@{}}
\toprule
\textbf{Q\#} & \textbf{Topic} & \textbf{Question} \\
\midrule
Q1 & e-Office Usage & Do people in AIGGPA use e-Office? \\
Q2 & Digitisation & What is the status of archiving and digitisation of files? \\
Q3 & Tools \& Efficiency & What kind of tech or tools do employees use? What are ways to increase their efficiency? \\
Q4 & Asset/File Mgmt & Is there an asset management tool or file management tool in use? \\
Q5 & Employee Needs & Do employees require any additional tools? (ref: Google Workspace) \\
Q6 & Tool Inventory & What other tools are the employees currently using? \\
\bottomrule
\end{tabularx}
\end{table}

% =============================================================================
\section{Scope \& Concept Proposal Direction}

\subsection{Target Sectors for Capacity Building Study}

The manager identified \textbf{7 sectors} for the capacity building solutions study:

\begin{enumerate}[label=\textbf{(\arabic*)}]
  \item Forest (\hindi{वन})
  \item Health (\hindi{स्वास्थ्य})
  \item Education (\hindi{शिक्षा})
  \item Agriculture (\hindi{कृषि})
  \item Rural Development (\hindi{ग्रामीण विकास})
  \item Urban Development (\hindi{शहरी विकास})
  \item State Office and District Office (\hindi{राज्य कार्यालय एवं जिला कार्यालय})
\end{enumerate}

\begin{actionbox}[Open Question from Manager]
\textbf{How many departments should we go focused on?}

Should the study cover officials across \textbf{Class IV, III, II, and I} -- or be narrowed?
\end{actionbox}

\subsection{Concept Proposal Structure (Manager's Direction)}

The manager outlined the following structure for the concept note/proposal:

\begin{enumerate}[label=\textbf{\arabic*.}]
  \item \textbf{Background} $\rightarrow$ Why? Purpose?
  \item \textbf{Government's Role:} Duties, governance obligations
  \item \textbf{Technology for Service Delivery:}
    \begin{itemize}
      \item Some app to send and receive files
      \item Look at how government can serve citizens better
      \item How government employees interact with schemes directly
    \end{itemize}
  \item \textbf{Current Technology Audit:}
    \begin{itemize}
      \item What tech is the government currently using -- and what do they \textit{think} about it?
    \end{itemize}
\end{enumerate}

\subsection{Two Dimensions of Technology Use}

The manager made a key distinction:

\begin{table}[h!]
\centering
\begin{tabularx}{\textwidth}{@{}lX@{}}
\toprule
\textbf{Dimension} & \textbf{Description} \\
\midrule
\textbf{Technology used to serve the public} & Platforms, portals, apps used for citizen-facing service delivery (G2C) \\
\textbf{Technology used for personal/internal use} & Tools used by employees for their own administrative work (G2G) \\
\bottomrule
\end{tabularx}
\end{table}

\subsection{Research Methodology Direction}

\begin{insightbox}[Manager's Guidance on Methodology]
\textbf{How to do the research?}

$\rightarrow$ \textbf{Primary research} -- \textbf{Exploratory} approach

This confirms the study should be field-based, qualitative/mixed-methods, using surveys, interviews, and direct observation at AIGGPA and target departments.
\end{insightbox}

% =============================================================================
\section*{Summary: Key Takeaways for Next Steps}

\begin{enumerate}[label=\textbf{\arabic*.}]
  \item \textbf{No existing research or data} on HR-IT integration at AIGGPA -- this is greenfield work
  \item \textbf{Primary exploratory research} is the confirmed methodology
  \item \textbf{3-tier employee classification:} Active Users / Aware Non-Users / Unaware
  \item \textbf{7 target sectors} identified for capacity building study
  \item \textbf{6 specific research questions} from the manager to guide survey design
  \item \textbf{Two technology dimensions:} public-facing (G2C) vs. internal (G2G)
  \item \textbf{Concept proposal} needed with Background, Purpose, Tech Audit, and Recommendations
  \item \textbf{Mission Karmayogi's 6 pillars} provide the national policy framework to anchor the study
\end{enumerate}

\end{document}
